% !TITLE 念奴娇·赤壁怀古
% !AUTHOR 苏轼
% !DYNASTY 两宋
% !GENRE 词
% !ORDER 7

\begin{poem}
大江东去，浪淘尽，千古风流人物\textsuperscript{1}。
故垒\textsuperscript{2}西边，人道是，三国周郎赤壁\textsuperscript{3}。
乱石穿空，惊涛拍岸，卷起千堆雪\textsuperscript{4}。
江山如画，一时多少豪杰。

遥想公瑾\textsuperscript{5}当年，小乔\textsuperscript{6}初嫁了，雄姿英发。
羽扇纶巾\textsuperscript{7}，谈笑间，樯橹\textsuperscript{8}灰飞烟灭。
故国神游\textsuperscript{9}，多情应笑我，早生华发\textsuperscript{10}。
人生如梦，一尊还酹\textsuperscript{11}江月。
\end{poem}

\begin{annotations}
\notetext{1}{风流人物：英雄人物，才华出众的人物。大江日夜冲刷，淘走了千古英雄的痕迹。}
\notetext{2}{故垒：古老的营垒，过去战争留下的壁垒。}
\notetext{3}{赤壁：三国时期赤壁之战的古战场，周瑜在此大败曹操。苏轼写此词时身处黄州（今湖北黄冈）的赤鼻矶，并非真正的赤壁古战场，但他借景怀古。}
\notetext{4}{千堆雪：无数堆白色的浪花。把浪花比作雪，写出江水的壮阔。}
\notetext{5}{公瑾：周瑜的字。周瑜（175-210），三国吴国名将，赤壁之战的总指挥。}
\notetext{6}{小乔：周瑜的妻子，与姐姐大乔并称"二乔"，是当时著名的美女。}
\notetext{7}{羽扇纶巾：手持羽毛扇、头戴青丝头巾。纶（guān），青色的丝带。这是儒将的潇洒打扮——周瑜不是赳赳武夫，而是文质彬彬的统帅。}
\notetext{8}{樯橹：战船的桅杆和船桨，代指曹操的战船。樯（qiáng），桅杆；橹（lǔ），大桨。周瑜用火攻，曹操的战船烧成灰烬。}
\notetext{9}{故国神游：心神游历到当年的古战场。}
\notetext{10}{华发：花白的头发。多情的人应该笑话我，这么早就白了头发——对英雄的遥想映照出自我的感慨。}
\notetext{11}{酹：把酒洒在地上，表示祭奠。酹（lèi）。举起酒杯，把酒洒向江月，祭奠那些消逝的英雄。}
\end{annotations}
